\documentclass[12pt,a4paper]{report}

\usepackage[T1]{fontenc}
\usepackage[utf8]{inputenc}
\usepackage{lmodern}
\usepackage[top=1in,bottom=1in,left=1.25in,right=1in]{geometry}
\usepackage{setspace}
\usepackage{graphicx}
\graphicspath{{figures/}}
\usepackage{booktabs}
\usepackage{array}
\usepackage{tabularx}
\usepackage{longtable}
\usepackage{multirow}
\usepackage{amsmath}
\usepackage{amssymb}
\usepackage{xcolor}
\usepackage{listings}
\lstset{
  basicstyle=\ttfamily\small,
  breaklines=true,
  frame=single,
  numbers=left,
  numberstyle=\tiny\color{gray},
  keywordstyle=\bfseries\color{blue!70!black},
  commentstyle=\itshape\color{green!50!black},
  stringstyle=\color{red!60!black},
  showstringspaces=false,
  captionpos=b
}
\usepackage[hidelinks]{hyperref}
\usepackage{fancyhdr}
\usepackage{titlesec}
\usepackage{caption}
\usepackage{subcaption}
\usepackage{float}
\usepackage{tikz}
\usepackage{pgfplots}
\pgfplotsset{compat=1.18}
\usetikzlibrary{shapes.geometric,arrows.meta,positioning,fit,backgrounds,calc}
\usepackage[style=ieee,backend=biber]{biblatex}
\addbibresource{references.bib}

\pagestyle{fancy}
\fancyhf{}
\fancyhead[L]{\small LeBlanc: Automated Prompt Enrichment and Iterative Repair}
\fancyhead[R]{\small \thepage}
\renewcommand{\headrulewidth}{0.4pt}

\titleformat{\chapter}[block]
  {\large\bfseries\centering}{Chapter \thechapter:}{0.5em}{}
\titlespacing*{\chapter}{0pt}{20pt}{10pt}

\onehalfspacing

% =========================================================
\begin{document}
\setcounter{page}{8}

% =========================================================
% CHAPTER 1 -- INTRODUCTION
% =========================================================
\chapter{Introduction}
\label{ch:introduction}

\section{Background and Motivation}
\label{sec:background}

The use of large language models as coding assistants has become standard
practice in professional software development. Tools such as GitHub Copilot,
ChatGPT, and Google Gemini generate syntactically correct, functionally plausible
code within seconds of receiving a natural-language prompt. Developer surveys
consistently indicate that a majority of professional programmers now use at
least one AI-assisted coding tool regularly, and this adoption is accelerating.

This adoption creates a security problem that is less visible but well-documented.
LLMs are trained primarily to maximise functional correctness---to produce code
that runs and satisfies the stated requirement. They learn from public repositories
on platforms such as GitHub, where a substantial fraction of existing code contains
known security weaknesses. The model learns these patterns alongside everything
else and reproduces them when asked to write similar code.

The consequence in practice is direct. A developer who asks an LLM to ``write a
Flask login endpoint with MySQL'' receives working code that is, with measurable
probability, also exploitable. SQL query parameters go unsanitised. Password
comparison happens in plaintext. Session tokens are returned without appropriate
flags. The developer sees the code run correctly against a test database, ships
it, and the vulnerability ships with it.

The 2022 study by Pearce et al.~\cite{pearce2022asleep} quantified this problem
against GitHub Copilot: approximately 40\% of the generated code samples they
evaluated contained at least one security vulnerability, with a vulnerability
rate of 63.4\% across the LLMSecEval prompt set. Subsequent studies on ChatGPT
and other models found comparable or higher rates~\cite{khoury2023secure,
tony2023llmseceval}. This is not a problem specific to one model or one point in
time---it appears to be a structural consequence of training on general-purpose
code repositories without security-specific supervision.

\section{Problem Statement}
\label{sec:problem}

Three partial solutions exist in the literature, each addressing one aspect of
the problem but not the full cycle.

The first class modifies the training process: fine-tuning on secure code
pairs~\cite{hexacoder2024}, preference optimisation toward secure
outputs~\cite{lpo2024}, or constrained decoding at inference
time~\cite{codeguard2024}. These approaches can be effective but require access
to model weights or specialised training infrastructure. They cannot be applied
to closed-source commercial models and require a new training run for each
model version.

The second class is evaluation-only: run a set of prompts through one or more
models, scan the outputs, and report a vulnerability rate. The
SecurityEval~\cite{securityeval}, LLMSecEval~\cite{tony2023llmseceval}, and
CodeSecEval~\cite{codeseceval2023} benchmarks all fall into this category. These
studies produce useful comparative data but stop at detection. They measure
vulnerability rates; they do not reduce them.

The third class addresses the prompt side: inject CWE-specific warnings before
sending the prompt to the model. The SecuCoGen study~\cite{tihanyi2023secucogenicse}
demonstrated this approach on GPT-3.5 and found substantial reduction. However,
the study used a single model, a fixed dataset of 180 samples, and included no
post-generation repair mechanism.

The gap all three classes leave open is: a single deployable system that
(a) enriches prompts with CWE warnings automatically, (b) scans generated code
with real static analysis tools, and (c) drives an iterative repair loop until
the scanner returns clean, running across multiple models simultaneously for
direct comparison. No such system exists in the published literature.

\section{Proposed System: LeBlanc}
\label{sec:proposed}

LeBlanc is a three-engine pipeline that fills this gap. The name references the
Le Blanc cost-of-quality principle: defects caught at the point of creation are
orders of magnitude cheaper to fix than defects found after deployment. The same
principle applies directly to security vulnerabilities in generated code---the
cost of not catching a vulnerability during generation is borne entirely by the
downstream developer and ultimately the end user.

The three engines are:

\begin{itemize}
  \item \textbf{Engine A --- Prompt Enrichment.} A rule-based local system that
        parses any coding prompt for security-relevant keywords and injects
        structured CWE warnings before the LLM call. The enrichment is
        deterministic and reproducible; no LLM is involved.

  \item \textbf{Engine B --- Static Analysis.} Once the LLM returns code,
        Engine~B runs Semgrep and Bandit via subprocess calls and maps their raw
        JSON output to structured findings: CWE identifier, severity, line number,
        and vulnerability category.

  \item \textbf{Engine C --- Iterative Repair.} When Engine~B finds
        vulnerabilities, Engine~C builds a repair prompt from the scanner findings
        and resubmits it to the originating LLM. The patched output is immediately
        re-scanned. The loop runs for up to three iterations or until the scanner
        returns clean.
\end{itemize}

The pipeline runs against multiple LLMs simultaneously. For any given prompt,
results for all models and all modes---plain, enriched, and enriched with
repair---are collected, scanned, and stored in a single database run. The
frontend provides a side-by-side comparison of vulnerability rates, scan
findings, and repair convergence for each model.

\section{Research Questions}
\label{sec:rqs}

LeBlanc is designed to answer five research questions from a single automated
data collection run:

\begin{description}
  \item[RQ1 (Enrichment Effect).] Does automated CWE-aware prompt enrichment
        reduce vulnerability rates, and by how much per model?
  \item[RQ2 (Model Comparison).] Which model produces the safest code out of
        the box? Does the ranking change after enrichment?
  \item[RQ3 (Repair Effectiveness).] How many repair iterations does each model
        need to reach clean code? Is there a ceiling effect where some vulnerability
        types do not converge?
  \item[RQ4 (CWE Difficulty).] Which CWE categories are hardest for LLMs to
        avoid and hardest to self-repair?
  \item[RQ5 (Generational Gap).] How does vulnerability rate and repair
        convergence differ between a small older model (Llama 3.1-8B) and
        contemporary frontier models (Command-A, Gemini 2.5 Flash)?
\end{description}

\section{Scope and Limitations}
\label{sec:scope}

The current evaluation covers 33 prompts from the LLMSecEval benchmark, three
models, and up to three repair iterations per run. The dataset and model set are
intentionally bounded to allow one complete reproducible run within semester
constraints. The LLMSecEval prompts are Python-only; Engine~B is correspondingly
configured with Python-specific static analysis rules.

The Llama 3.1-8B results require careful interpretation. An 8-billion parameter
model represents an earlier, smaller generation of open-weight LLMs. Its behaviour
under prompt enrichment---where adding CWE warnings increased rather than reduced
vulnerability rate---is consistent with published evidence that small models
degrade under complex instruction sets rather than improving. These results are
included for completeness and as a motivating datapoint for RQ5. They should not
be taken as evidence that prompt enrichment is counterproductive in general.
Future work will repeat the evaluation on Llama 3.3-70B, which uses the same
Llama architecture family at a scale known to follow complex instructions reliably.

\section{Contributions}
\label{sec:contributions}

This project makes three concrete contributions:

\begin{enumerate}
  \item \textbf{A working pipeline system.} LeBlanc is a fully implemented,
        runnable web application that takes a coding prompt, runs it through all
        three engines across multiple LLMs in parallel, and displays structured
        results. The source code and database schema are reproducible from the
        submitted repository.

  \item \textbf{An empirical dataset.} Every pipeline run logs the original
        prompt, enriched prompt, raw LLM output, extracted code, Semgrep and
        Bandit findings, and every repair iteration to a SQLite database. The
        collected dataset covers 321 total runs across three models, three modes,
        and 33 prompts.

  \item \textbf{Comparative results across eight metrics.} Vulnerability Rate,
        Relative Improvement over the Copilot baseline, Prompt Enrichment Gain,
        Convergence Rate at one and three iterations, Mean Iterations to Fix,
        Remaining Failure Rate, Agreement Rate, and Unique Failure Rate are all
        reported per model and per CWE category.
\end{enumerate}

\section{Report Organisation}
\label{sec:organisation}

Chapter~\ref{ch:literature} reviews the eight prior works that directly motivate
LeBlanc's design and maps each to the specific gap it leaves.
Chapter~\ref{ch:system} describes the system architecture, the three engines, and
implementation decisions.
Chapter~\ref{ch:methodology} defines the evaluation methodology: dataset
selection, metrics definitions, and experimental setup.
Chapter~\ref{ch:results} presents results for all five research questions.
Chapter~\ref{ch:discussion} interprets the findings, discusses limitations, and
addresses the Llama 3.1-8B result specifically.
Chapter~\ref{ch:conclusion} summarises the work and outlines future directions
including the VS~Code extension and the planned larger-model evaluation.

% =========================================================
% CHAPTER 2 -- LITERATURE REVIEW
% =========================================================
\chapter{Literature Review}
\label{ch:literature}

This chapter reviews the eight papers that directly informed LeBlanc's design.
Each is summarised on its contribution and the specific gap it leaves open.
The gaps are not criticisms---each paper is well-executed within its stated scope.
They define the boundary of existing knowledge that LeBlanc extends.

\section{Overview}
\label{sec:litoverview}

The review spans three themes: (1) vulnerability measurement studies that
establish baseline rates, (2) enrichment and fine-tuning approaches that attempt
to reduce rates at generation time, and (3) repair and feedback approaches that
attempt to fix vulnerabilities after generation. LeBlanc operates across all
three themes simultaneously.

Table~\ref{tab:litgap} summarises each paper against the capabilities present
in LeBlanc.

\begin{table}[H]
\centering
\caption{Prior work and the gaps addressed by LeBlanc}
\label{tab:litgap}
\small
\renewcommand{\arraystretch}{1.3}
\begin{tabularx}{\textwidth}{@{}p{1.6in}p{1.5in}p{1.5in}p{1.5in}@{}}
\toprule
\textbf{Paper} & \textbf{Core Contribution} & \textbf{Gap Left} & \textbf{LeBlanc's Answer} \\
\midrule
Pearce et al.\ (2022) \cite{pearce2022asleep}
  & Copilot audit; 63.4\% VR baseline
  & Single model; no repair; observational
  & Multi-model; automated repair; active pipeline \\

Tony et al.\ (2023) \cite{tony2023llmseceval}
  & LLMSecEval dataset; 120 NL prompts
  & Single-turn; no enrichment or repair
  & Enrichment + repair on LLMSecEval prompts \\

Khoury et al.\ (2023) \cite{khoury2023secure}
  & ChatGPT across 21 scenarios
  & Single model; manual revision only
  & Multi-model; automated repair loop \\

Tihanyi et al.\ (2024) \cite{tihanyi2023secucogenicse}
  & CWE-aware prompting; VR 31\% $\to$ 14\%
  & Single model; 180 samples; no repair
  & Enrichment automated; multi-model; + repair \\

Wang et al.\ (2024) \cite{codeseceval2023}
  & 44-CWE benchmark; automated evaluation
  & No enrichment; no iterative repair
  & All three modes in one system \\

He \& Vechev (2024) \cite{codeguard2024}
  & Constrained decoding at inference time
  & Requires model internals; no API support
  & API-only; works on any black-box model \\

Siddiq et al.\ (2024) \cite{hexacoder2024}
  & Oracle-guided secure fine-tuning
  & Fine-tune only; closed models excluded
  & No fine-tuning; prompt-side only \\

Zhou et al.\ (2024) \cite{lpo2024}
  & Preference optimisation for security
  & Requires retraining; not deployable
  & Fully deployment-agnostic \\
\bottomrule
\end{tabularx}
\end{table}

\section{Vulnerability Measurement Studies}
\label{sec:litbaseline}

\subsection{Pearce et al.\ (2022) --- Asleep at the Keyboard}
\label{ssec:pearce}

Pearce et al.~\cite{pearce2022asleep} conducted the first systematic security
audit of GitHub Copilot's code generation output. They constructed 89 scenarios
spanning CWE categories including injection, memory safety, and cryptographic
weaknesses and evaluated Copilot's outputs through manual expert review and
automated scanning. The study found that approximately 40\% of all generated
code samples contained at least one security vulnerability, and that certain
categories---particularly CWE-89 (SQL Injection) and CWE-22 (Path
Traversal)---appeared at high rates across multiple scenarios.

The headline figure from the LLMSecEval subset is a vulnerability rate of 63.4\%:
across 82 prompts, 52 produced vulnerable output. This number is the standard
external baseline for subsequent LLM security evaluation work. LeBlanc uses it
as the primary comparison point throughout the results chapter, and the
metrics engine computes Relative Improvement (RI) against this value for every
model and mode combination.

The study's limitation from LeBlanc's perspective is that it is observational
and single-model. It measures vulnerability rates but proposes no mechanism to
reduce them. There is no enrichment step, no repair loop, and no comparison
between models. It establishes the scale of the problem but does not address it.

\subsection{Tony et al.\ (2023) --- LLMSecEval}
\label{ssec:tony}

Tony et al.~\cite{tony2023llmseceval} produced the LLMSecEval dataset: 120
natural-language prompts designed to elicit code containing common security
weaknesses from the CWE Top 25. Each prompt is annotated with its target CWE
and a ``vulnerable by design'' flag indicating that the expected model output,
without security guidance, will contain the weakness. The dataset covers Python
scenarios for weaknesses including SQL injection, path traversal, hardcoded
credentials, and weak cryptography.

LeBlanc uses a 33-prompt subset of LLMSecEval as its primary evaluation dataset.
The prompts are functional application requests---``write a Flask login
endpoint,'' ``build a file upload handler''---rather than explicitly asking the
model to produce insecure code, which would risk triggering safety filters in
commercial models and make the results non-representative of actual developer
usage.

Tony et al.'s own evaluation is single-turn: each prompt is sent to several
models and the output is classified as containing or not containing the target
CWE. No enrichment, repair, or multi-turn analysis is performed. LeBlanc extends
this by running each prompt through all three pipeline modes and logging the full
scan trace for each run.

\subsection{Khoury et al.\ (2023) --- How Secure Is Code Generated by ChatGPT?}
\label{ssec:khoury}

Khoury et al.~\cite{khoury2023secure} evaluated ChatGPT across 21 scenarios
covering buffer overflows, SQL injection, cryptographic weaknesses, and
other common CWE types. They found that approximately 16 of the 21 scenarios
produced at least one piece of code containing a security vulnerability on the
first request. Notably, they also found that explicitly requesting a revision---``make
this code secure''---caused ChatGPT to produce a safer version in several cases.

The revision step is manual and qualitative; no scanner is used to verify whether
the revision actually removes the vulnerability, and no iteration count is
tracked. This observation nonetheless provides empirical support for the core
premise of Engine~C: the same LLM that produces vulnerable code can, when given
specific feedback about what is wrong, produce a corrected version. LeBlanc
automates and verifies this process.

\section{Enrichment and Prompt-Side Approaches}
\label{sec:litenrich}

\subsection{Tihanyi et al.\ (2024) --- SecuCoGen}
\label{ssec:secucogen}

Tihanyi et al.~\cite{tihanyi2023secucogenicse} is the most directly relevant
prior work to LeBlanc's Engine~A. They constructed a 180-sample Python dataset
of security-sensitive prompts and evaluated GPT-3.5-Turbo under two conditions:
standard prompts, and CWE-enriched prompts that explicitly named the vulnerability
the model should avoid. Under standard prompts, GPT-3.5 produced vulnerable code
31\% of the time. Under CWE-enriched prompts this dropped to 14\%---a reduction
of 55\%.

This result is the primary empirical justification for Engine~A's design. If
telling the model which vulnerability to avoid in the prompt reduces vulnerability
rates by over half on one model, the natural next question is whether this effect
generalises across models and whether it can be automated for arbitrary input
prompts.

LeBlanc automates the SecuCoGen enrichment strategy through a keyword-to-CWE
mapping file: Engine~A parses any incoming prompt for security-relevant keywords
and injects structured warnings automatically, without requiring the prompt
author to know which CWEs are relevant. SecuCoGen's limitations---single model,
fixed Python-only dataset, no repair---are all directly addressed.

\subsection{Siddiq et al.\ (2024) --- HexaCoder}
\label{ssec:hexacoder}

HexaCoder~\cite{hexacoder2024} takes a training-time approach to reducing
vulnerability rates. The system generates a large training set of (vulnerable,
secure) code pairs using an ``oracle''---a combination of scanner feedback and
expert review---and fine-tunes a base LLM on this dataset. The resulting model
generates substantially safer code than the untuned baseline.

The critical limitation for LeBlanc's use case is deployment: fine-tuning
requires access to model weights. It cannot be applied to closed-source commercial
APIs such as Gemini or Command-A, which are the most widely used models in
practice. It also requires substantial compute resources and a retraining run
for every new model version. LeBlanc rejects any fine-tuning dependency: all
three engines operate at the API level and work against any model accessible via
a standard chat completion endpoint.

\subsection{Zhou et al.\ (2024) --- LPO}
\label{ssec:lpo}

Zhou et al.~\cite{lpo2024} propose Language Model Preference Optimisation (LPO),
a variant of reinforcement learning from human feedback applied to the security
domain. A reward model scores generated code on security criteria, and the base
LLM is fine-tuned to maximise that reward. The approach produces models with
improved security properties compared to standard instruction-tuned baselines.

The same deployment limitation applies as for HexaCoder. LPO is an effective
research contribution for scenarios where model weights are available, but it
provides no mechanism for improving security on the closed commercial models
that most developers actually use.

\section{Post-Generation and Constrained Approaches}
\label{sec:litpost}

\subsection{He and Vechev (2024) --- CODEGUARD+}
\label{ssec:codeguard}

He and Vechev~\cite{codeguard2024} present CODEGUARD+, a constrained decoding
approach that modifies the token sampling process during generation to steer the
model away from patterns associated with known vulnerabilities. The constraints
are derived from a learned security classifier that operates on partial token
sequences during decoding.

CODEGUARD+ is technically sophisticated and measurably effective within its
evaluation. However, it requires direct access to the model's logit distribution
at each decoding step, which is only possible when running the model locally.
This rules out all API-served commercial models. LeBlanc operates entirely at
the output level---it reads generated text, not intermediate model states---and
therefore has no such constraint.

\subsection{Wang et al.\ (2024) --- CodeSecEval}
\label{ssec:codeseceval}

Wang et al.~\cite{codeseceval2023} developed CodeSecEval, a benchmark covering
44 CWE types across 180 samples. The benchmark includes both a generation task
(produce code satisfying a description) and a repair task (fix a provided
vulnerable snippet). They evaluate several models on both tasks and report
per-CWE difficulty rankings.

CodeSecEval is the broadest existing benchmark by CWE coverage, and its per-CWE
difficulty rankings directly inform LeBlanc's RQ4 analysis. The key distinction
from LeBlanc's Engine~C is that CodeSecEval's repair task is static: a fixed
vulnerable snippet is provided to the model as input, and one repair attempt is
recorded. There is no iterative loop, no re-scanning after the attempt, and no
automated verification that the repair succeeded. Engine~C closes this gap by
re-scanning after every repair attempt and continuing until the scanner confirms
clean output or the iteration limit is reached.

\section{Summary and Positioning}
\label{sec:litgapsummary}

Five consistent gaps emerge from this review:

\begin{enumerate}
  \item \textbf{Single-model evaluation.} Six of the eight papers evaluate one
        model. No prior work provides cross-model comparative data under identical
        conditions.
  \item \textbf{Single-turn evaluation.} Seven of the eight papers test generation
        once and stop. No study measures iterative automated repair with
        re-scanning.
  \item \textbf{No automated enrichment.} SecuCoGen shows that CWE-aware prompting
        works but requires manual prompt construction. No prior system automates
        enrichment from an arbitrary input prompt.
  \item \textbf{No integrated system.} Enrichment, scanning, and repair exist in
        separate papers; they have not been assembled into a single runnable tool.
  \item \textbf{No logged multi-model dataset.} No prior study logs all
        runs---prompts, enriched prompts, raw outputs, scan findings, repair
        iterations---to a queryable database.
\end{enumerate}

LeBlanc addresses all five. The combination of automated enrichment across
multiple models, real static analysis, iterative repair with re-scanning, a
deployable web application, and a logged reproducible dataset has not been
presented in any prior work.

\printbibliography[heading=bibintoc,title={References}]

\end{document}
